\documentclass[11pt,a4paper]{article}

\usepackage[utf8]{inputenc}
\usepackage[T1]{fontenc}
\usepackage{lmodern}
\usepackage[margin=0.72in]{geometry}
\usepackage[table]{xcolor}
\usepackage{array}
\usepackage{tabularx}
\usepackage{booktabs}
\usepackage{enumitem}
\usepackage{titlesec}
\usepackage{fancyhdr}
\usepackage{lastpage}
\usepackage{hyperref}

\definecolor{TextInk}{HTML}{233142}
\definecolor{HeadingBlue}{HTML}{1D3557}
\definecolor{AccentTeal}{HTML}{2A7F92}
\definecolor{AccentGold}{HTML}{C58B4E}
\definecolor{RuleGray}{HTML}{D6E0E5}
\definecolor{TableStripe}{HTML}{F3F8FA}
\definecolor{SoftTint}{HTML}{F7F4EE}
\definecolor{SoftBlue}{HTML}{EEF6F8}
\definecolor{SoftGreen}{HTML}{EEF7F0}
\definecolor{SoftRed}{HTML}{FBF0EF}

\hypersetup{
  colorlinks=true,
  linkcolor=HeadingBlue,
  urlcolor=AccentTeal,
  pdftitle={IB Geography SL May 4 Remediation Workbook},
  pdfauthor={IB Geography SL Preparation}
}

\color{TextInk}
\setlength{\parskip}{0.42em}
\setlength{\parindent}{0pt}
\setlength{\tabcolsep}{5pt}
\renewcommand{\arraystretch}{1.16}
\setlist[itemize]{leftmargin=1.32em,itemsep=3pt,topsep=4pt}
\setlist[enumerate]{leftmargin=1.45em,itemsep=4pt,topsep=5pt}

\newcolumntype{Y}{>{\raggedright\arraybackslash}X}
\newcolumntype{L}[1]{>{\raggedright\arraybackslash}p{#1}}
\newcolumntype{C}[1]{>{\centering\arraybackslash}p{#1}}

\titleformat{\section}
  {\normalfont\Large\sffamily\bfseries\color{HeadingBlue}}
  {}{0pt}{}
  [\vspace{0.08em}{\color{AccentGold}\titlerule[0.8pt]}]

\titleformat{\subsection}
  {\normalfont\large\sffamily\bfseries\color{AccentTeal}}
  {}{0pt}{}

\pagestyle{fancy}
\setlength{\headheight}{14pt}
\fancyhf{}
\fancyhead[L]{\sffamily\color{AccentTeal}May 4 Workbook}
\fancyhead[R]{\sffamily\color{HeadingBlue}IB Geography SL}
\fancyfoot[C]{\sffamily\color{HeadingBlue}\thepage\ of \pageref{LastPage}}
\renewcommand{\headrulewidth}{0.4pt}

\newcommand{\term}[1]{\textcolor{HeadingBlue}{\textbf{#1}}}
\newcommand{\exam}[1]{\textcolor{AccentTeal}{\textbf{#1}}}
\newcommand{\todobox}{\fbox{\rule{0pt}{0.8em}\rule{0.8em}{0pt}}}
\newcommand{\smallline}{\rule{0.96\linewidth}{0.4pt}}
\newcommand{\blankbox}[2]{%
\noindent\fcolorbox{RuleGray}{#1}{%
  \begin{minipage}[t][#2][t]{0.965\textwidth}
  \vspace{0.2em}
  \end{minipage}
}}
\newcommand{\boxnote}[3]{%
\noindent\colorbox{#1}{%
  \begin{minipage}{0.965\textwidth}
  \sffamily
  \textbf{\textcolor{HeadingBlue}{#2}} #3
  \end{minipage}
}}

\begin{document}

\begin{center}
  {\sffamily\bfseries\color{HeadingBlue}\LARGE May 4 Remediation Workbook}\\[0.25em]
  {\sffamily\color{AccentTeal}\large Urban Environments Redo + Population Data Drill}
\end{center}

\boxnote{SoftBlue}{Use with the May 4 study pack.}{For the strict
\term{25 minute} block, choose \term{one} main task: Task A, Task B, or Task C.
If you are following the May 2 teacher feedback, choose Task B unless Q1 source
precision was clearly the biggest weakness. Task A can be used as a short
warm-up or extra-time follow-up outside the timed block.}

\section{Choose One Weakest Question Type}

\rowcolors{2}{TableStripe}{white}
\begin{tabularx}{\textwidth}{L{0.11\textwidth}Y L{0.24\textwidth}}
\toprule
\textbf{Use} & \textbf{Weakness} & \textbf{Task} \\
\midrule
\todobox & I lost short-answer marks because I gave a general idea without
exact source evidence. & Task A \\
\todobox & I lost 10-mark marks because my Detroit evidence or evaluation was
too general. & Task B \\
\todobox & I lost Paper 2 marks because population data, models, or comparisons
were weak. & Task C \\
\bottomrule
\end{tabularx}
\rowcolors{2}{}{}

\section{Task A: Source Evidence Redo}

\subsection{Source A: Urban Environmental Quality}

The data below are illustrative and are designed to practise the May 2 weak
question type.

\rowcolors{2}{TableStripe}{white}
\begin{tabularx}{\textwidth}{L{0.27\textwidth}C{0.15\textwidth}C{0.15\textwidth}C{0.15\textwidth}C{0.15\textwidth}}
\toprule
\textbf{Indicator} & \textbf{CBD} & \textbf{Inner-city regeneration zone} &
\textbf{Suburban zone} & \textbf{Peripheral informal settlement} \\
\midrule
Population density index & 88 & 72 & 34 & 96 \\
Green space share & 8\% & 14\% & 34\% & 3\% \\
Rent index & 91 & 82 & 67 & 24 \\
Service access & High & Medium-high & High & Weak \\
Main environmental pressure & Congestion and limited open space &
Redevelopment pressure & Car dependence and land take & Flood / pollution
exposure \\
\bottomrule
\end{tabularx}
\rowcolors{2}{}{}

\begin{enumerate}[label=\textbf{A\arabic*.}]
  \item \exam{State} two pieces of evidence from Source A that show uneven
  environmental quality. \hfill \textbf{[2]}
\end{enumerate}

\blankbox{white}{2.5cm}

\subsection{Self-Check}

\begin{tabularx}{\textwidth}{L{0.12\textwidth}Y}
\toprule
\textbf{Done} & \textbf{Check} \\
\midrule
\todobox & I named a specific zone or category. \\
\todobox & I used an exact number or exact source phrase. \\
\todobox & I made the unevenness clear by comparison. \\
\bottomrule
\end{tabularx}

\section{Task B: Detroit 10-Mark Redo}

\subsection{Prompt}

\boxnote{SoftTint}{Paper 1 extended-answer practice.}{\exam{Evaluate} the
extent to which regeneration or gentrification can solve problems caused by
urban decline. Use one or more named urban areas. \hfill \textbf{[10]}}

\subsection{Plan First: 4 Minutes}

\begin{tabularx}{\textwidth}{L{0.24\textwidth}Y}
\toprule
\textbf{Need} & \textbf{Plan notes} \\
\midrule
Command term and judgement & \smallline \\[0.7em]
Detroit decline evidence & \smallline \\[0.7em]
Regeneration benefit & \smallline \\[0.7em]
Gentrification / uneven-impact limit & \smallline \\[0.7em]
Final judgement phrase & \smallline \\[0.7em]
\bottomrule
\end{tabularx}

\subsection{Write: 12 Minutes}

\blankbox{white}{13.2cm}

\newpage

\subsection{Paragraph Upgrade Checklist}

\rowcolors{2}{TableStripe}{white}
\begin{tabularx}{\textwidth}{L{0.12\textwidth}Y}
\toprule
\textbf{Done} & \textbf{Check} \\
\midrule
\todobox & I named Detroit, Michigan, United States. \\
\todobox & I used at least one precise evidence anchor: \term{1.85 million in
1950}, \term{639,111 in 2020}, \term{32.7\% poverty}, or another verified
class statistic. \\
\todobox & I explained decline through deindustrialization, suburbanization,
population loss, vacancy, or fiscal stress. \\
\todobox & I explained a benefit of regeneration or gentrification. \\
\todobox & I evaluated a limitation: affordability, displacement, exclusion,
uneven benefits, or remaining poverty / service gaps. \\
\todobox & My conclusion uses degree language: \term{partly}, \term{only to
some extent}, \term{more successful for X than Y}. \\
\bottomrule
\end{tabularx}
\rowcolors{2}{}{}

\section{Task C: Population Data Drill}

\subsection{Source B: Population Indicators}

The data below are illustrative and are designed to practise Paper 2 comparison
skills.

\rowcolors{2}{TableStripe}{white}
\begin{tabularx}{\textwidth}{L{0.25\textwidth}C{0.17\textwidth}C{0.17\textwidth}C{0.17\textwidth}}
\toprule
\textbf{Indicator} & \textbf{Country Alpha} & \textbf{Country Beta} &
\textbf{Country Gamma} \\
\midrule
Total fertility rate & 4.6 & 1.3 & 2.1 \\
Median age & 18 & 48 & 32 \\
Annual population growth & 2.7\% & -0.3\% & 0.9\% \\
Net migration rate & -1.2\% & 0.4\% & 0.8\% \\
Population aged under 15 & 43\% & 12\% & 24\% \\
Population aged 65 and over & 3\% & 29\% & 12\% \\
\bottomrule
\end{tabularx}
\rowcolors{2}{}{}

\begin{enumerate}[label=\textbf{C\arabic*.}]
  \item Using Source B, \exam{describe} two contrasts between Country Alpha and
  Country Beta. \hfill \textbf{[4]}
  \item \exam{Explain} one reason why Country Gamma may continue to grow even
  though its fertility rate is close to replacement level. \hfill \textbf{[2]}
\end{enumerate}

\blankbox{white}{6.3cm}

\section{Final Error-Log Sentence}

Write one sentence you will reuse on exam day.

\boxnote{SoftGreen}{Examples.}{For short source questions, I will name the
category, indicator, and number. For Detroit essays, I will use population
decline plus uneven regeneration before judging the extent of success. For
Population data, I will compare with exact indicators before explaining.}

\blankbox{white}{2.8cm}

\end{document}
